\section{Methodology}
\label{sec:methods}

\subsection{CARLA scenario and system architecture}

Experiments use a controlled give-way junction in CARLA's Town05
\citep{dosovitskiy2017carla}. The ego follows a fixed route and encounters an
assertive or reactive target vehicle. Fixed-risk and adaptive-risk policies
share the trajectory solver and safety supervisor, making the executed motion
a coupled system outcome. \TODO{Insert CARLA version, synchronous step,
vehicle blueprints, route coordinates, weather, solver settings, supervisor
thresholds and A3 authority-regime definition from the frozen configuration.}

\subsection{Data collection and representation}

The factorial collection contains 200 rollouts: 50 ego initialisations crossed
with two target styles and two collection policies. Each prediction sample
contains a deployment-equivalent raster and six interaction tokens covering
$-1.0$ to $0.0$~s. Each token has 12 relative position, velocity, yaw and
distance features. The target future covers ten steps at 0.2~s intervals, for
a 2.0~s horizon, in the current target-local right-handed frame. Target style,
policy, factorial cell, split and initialisation ID are prohibited inputs.

Rollouts, not windows, are split: initialisations 1--40 train, 41--45
validation and 46--50 test. Normalisation is fit on train only. Samples with no
future label are excluded; partial-horizon samples use a masked loss.

\begin{figure}[t]
  \centering
  \includegraphics[width=0.98\linewidth]{figure01_research_workflow.png}
  \caption{Research workflow and evidence gates. The study progresses from
  the original planning question through controlled data, matched model
  comparison and frozen closed-loop evaluation; mixed and negative findings
  remain part of the evidence chain.}
  \label{fig:workflow}
\end{figure}

\begin{table}[t]
  \caption{Rollout-disjoint dataset split and usable prediction windows}
  \label{tab:dataset}
  \centering
  \begin{tabular}{lrrrrr}
    \toprule
    Split & Init IDs & Rollouts & Raw & Usable & Full horizon \\
    \midrule
    Train & 1--40 & 160 & 9,121 & 4,036 & 2,596 \\
    Validation & 41--45 & 20 & 1,034 & 506 & 326 \\
    Test & 46--50 & 20 & 1,075 & 495 & 315 \\
    \bottomrule
  \end{tabular}
\end{table}

\subsection{Prediction models}

\Bzero{} is the frozen pretrained MultiPath control. \Bone{} retrains only the
final MultiPath dense prediction head. \BtwoM{} and \BtwoD{} add flattened-
sequence MLP adapters; the former updates trajectory means and the latter
updates means, covariance/orientation and mixture logits. \Tone{} and \Ttwo{}
replace the corresponding MLP adapter with a masked Transformer. The
\BtwoM{}/\Tone{} and \BtwoD{}/\Ttwo{} parameter gaps are 10.48\% and 5.89\%,
below the preregistered 20\% matched-capacity ceiling.

\begin{figure}[t]
  \centering
  \includegraphics[width=0.98\linewidth]{figure02_model_architecture.png}
  \caption{Matched architecture and capacity controls. \Bone{} is the simple
  adaptation control; the MLP/Transformer pairs isolate attention from
  approximate adapter capacity while sharing the base prediction contract.}
  \label{fig:architecture}
\end{figure}

\TODO{Add an exact hyperparameter table: layer widths, heads, embedding size,
dropout, optimiser, learning rate, batch size, loss terms, epoch/early-stopping
rule, anchor count, frozen/trainable parameters and software versions.}

\subsection{Training, selection and calibration}

Five candidate variants are trained with seeds 11, 23 and 37 under the same
split, anchors, decoder, loss, optimiser budget and validation ranking metric.
The variant is selected by median validation rollout-macro NLL and a
representative seed is chosen using validation only. Selection is frozen before
the test set is opened. Context zero/shuffle ablations and collision-filtered
sensitivity runs are post-selection diagnostics and do not alter the primary
model selection.

Calibration parameters are fit on validation and frozen for test/deployment.
The closed-loop contrast uses the identity-calibrated \Bzero{} stack and the
validation-calibrated \Bone{} stack; it must therefore be interpreted as a
frozen predictor-stack treatment rather than a weights-only intervention.

\subsection{Statistical analysis}

Effect estimates can average the balanced paired conditions, but uncertainty
and exact sign-flip tests operate on five ego-initialisation cluster means.
The smallest attainable non-zero two-sided exact p-value is therefore 0.0625.
Holm correction is applied within each predefined family. Results emphasise
effect sizes, cluster intervals, sign consistency and exact/Holm p-values;
windows and simulator steps are never treated as independent replications.

\subsection{Data-quality audit and sensitivity}

Logging recorded 253 collision callbacks in six reactive training rollouts.
Episode reconstruction reduced these to 16 contact episodes and 62 unique
frames involving the target and infrastructure, not 253 ego--target crashes.
Because old callbacks cannot be aligned exactly to the sample clock, all 162
usable windows from the six rollouts are conservatively removed in a
validation-only sensitivity analysis. The primary dataset remains frozen; the
sensitivity asks only whether the architecture ranking changes.
